% §3 The extraction mapping tau
The extraction mapping $\tau$ converts a standard observability export into a grammar walk.
This section defines its three stages --- normalization, state inference, walk construction ---
as shipped in the open-source implementation,\footnote{\url{https://github.com/iamdgarcia/trace2grammar}}
and states the assumption that governs what the paper's coverage results inherit.
Figure~\ref{fig:pipeline} shows the pipeline.

\begin{figure}[t]
\centering
\resizebox{\textwidth}{!}{% The extraction pipeline: three observability dialects in, coverage atlas out.
% Layout: positioning-relative (border distances), so boxes cannot overlap no matter
% how wide their labels set them. Bottom row is a boustrophedon snake (right to left),
% which removes the old long curved arrow from features back to PCA.
\begin{tikzpicture}[>=Stealth, font=\small,
  box/.style={draw, rounded corners, align=center, minimum height=8mm},
  stage/.style={draw, rounded corners, align=center, minimum height=11mm, fill=black!4},
  ev/.style={->,>=Stealth},
  node distance=7mm and 9mm]
  % dialects (left, stacked)
  \node[box] (otel) {OTel GenAI\\spans};
  \node[box, below=7mm of otel] (oinf) {OpenInference\\spans};
  \node[box, below=7mm of oinf] (lsm) {LangSmith\\run trees};
  % top row, chained by borders (middle dialect anchors the entry stage)
  \node[stage, right=16mm of oinf] (norm) {normalize\\one mapper per dialect};
  \node[stage, right=of norm] (evts) {canonical events\\$\{t_{\mathrm{ms}},\ \mathrm{kind},\ \mathrm{name},\ \mathrm{attrs}\}$};
  \node[stage, right=of evts] (walk) {state walk $\sigma$\\$q_0 \to \cdots \to \theta$, outcome, $D$};
  \node[stage, right=of walk] (feat) {features\\grammar $|$ bigram};
  % bottom row (snake, right to left)
  \node[stage, below=11mm of feat] (proj) {PCA (2-D)\\fit on real set};
  \node[stage, below=11mm of walk] (grid) {density grid\\Gaussian bumps + blur};
  \node[stage, below=11mm of evts] (cov) {coverage \%\\+ atlas + top-$K$};
  % arrows
  \draw[ev] (otel.east) -- ++(4mm,0) |- (norm.west);
  \draw[ev] (oinf.east) -- (norm.west);
  \draw[ev] (lsm.east) -- ++(4mm,0) |- (norm.west);
  \draw[ev] (norm) -- (evts) node[midway,above,font=\scriptsize]{dedupe};
  \draw[ev] (evts) -- (walk) node[midway,above,font=\scriptsize]{rule classifier};
  \draw[ev] (walk) -- (feat) node[midway,above,font=\scriptsize]{encode};
  \draw[ev] (feat) -- (proj) node[midway,right,font=\scriptsize]{real-set basis};
  \draw[ev] (proj) -- (grid);
  \draw[ev] (grid) -- (cov);
\end{tikzpicture}
}
\caption{The extraction-and-coverage pipeline.
Three observability dialects (OTel GenAI, OpenInference, LangSmith; $\tau$-bench episodes as a
fourth evaluation-side dialect) are normalized into a canonical event stream, classified into a
grammar walk with outcome class and cycle depth, encoded as feature vectors, projected onto a
2-D plane fit on the real set, and scored for region coverage.}
\label{fig:pipeline}
\end{figure}

\subsection{Canonical event stream}
Every dialect is reduced to one format-agnostic event stream
$e = (t_{\mathrm{ms}}, \mathrm{kind}, \mathrm{name}, \mathrm{attrs})$ with event kinds
\begin{equation}
\begin{aligned}
\mathrm{kind} \in \{\,&\texttt{user\_msg},\ \texttt{assistant\_msg},\ \texttt{tool\_call},\\
&\texttt{tool\_result},\ \texttt{tool\_error},\ \texttt{hangup},\\
&\texttt{handoff},\ \texttt{session\_end}\,\}.
\end{aligned}
\label{eq:kinds}
\end{equation}
One normalization function per dialect maps its export shape onto
Equation~\eqref{eq:kinds}:
\begin{itemize}
  \item \texttt{\_normalize\_otel} --- OTel GenAI spans, conversation-scoped, with the
        \texttt{gen\_ai.*} attributes of \texttt{invoke\_agent} and
        \texttt{execute\_tool} calls;
  \item \texttt{\_normalize\_openinference} --- typed OpenInference spans:
        \texttt{input.value}, \texttt{output.value};
  \item \texttt{\_normalize\_langsmith} --- runs grouped by \texttt{parent\_run\_id};
  \item \texttt{\_normalize\_taubench} --- transcript turns plus the environment-verified
        reward, which drives the terminal outcome (a documented public analogue for
        production containment).
\end{itemize}
Timestamps are normalized to milliseconds (epoch nanoseconds or ISO timestamps depending on
dialect), attributes are flattened, and adjacent duplicate events are deduplicated.
Everything downstream of the canonical stream is format-agnostic; adding a dialect is one
normalization function, which is what keeps $\tau$ maintainable as the specification landscape
shifts --- the OTel GenAI conventions are still under Development status
\citep{otelgenai2026}.

\subsection{State inference as a config-exposed classifier}
The second stage maps the canonical stream to a grammar walk
$\sigma = q_0 \to q_1 \to \cdots \to \theta$ with states
$Q = \{q_{\mathrm{greet}}, q_{\mathrm{auth}}, q_{\mathrm{intent}}, q_{\mathrm{tool}},$
$q_{\mathrm{clarify}}, q_{\mathrm{confirm}}, q_{\mathrm{close}}\}$ and terminals
$\Theta = \{\theta_{\mathrm{cont}}, \theta_{\mathrm{aban}}, \theta_{\mathrm{hand}},$
$\theta_{\mathrm{esc}}\}$ \citep{garcia2026grammar}.
The seven states are not a domain model: they are the domain-independent phases of a
tool-calling dialogue (opening, identity, intent, execution, repair, confirmation, closing).
The set is a superset --- a walk uses the sub-sequence the deployment exhibits, so an agent
without authentication simply never emits $q_{\mathrm{auth}}$ --- and what varies per vertical
is which surface forms trigger each transition, never the phases themselves.
The baseline classifier is deliberately rule-based: keyword and regex pattern tables
(authentication tool names, handoff tool names, user clarification and confirmation phrases,
assistant closing phrases, explicit terminal event names) plus sequencing rules (a user
clarification re-opens the current settled state and increments cycle depth; tool errors with
recovery stay in $q_{\mathrm{tool}}$; explicit hangup and handoff events drive the terminal).
Every pattern lives in an \texttt{ExtractionConfig} dataclass, so a domain owner adapts the
mapping to a new vertical (banking, retail, IT support) by editing configuration, not code.

This is a heuristic, and the paper's claim does not depend on it being more.
The contribution is the \emph{interface}: deterministic walks with typed transitions that any
downstream measurement can consume.
The learned emission mapper that will replace the pattern tables is future work
(Section~\ref{sec:open_problems}); the configuration tables double as its training-data
generator.

\begin{definition}[Extraction mapping]
\label{def:tau}
$\tau : \mathcal{D} \to \Sigma^* \times \Theta \times \mathbb{N}$ maps an observability export
$d \in \mathcal{D}$ to a walk $\sigma$ over states $q \in Q$ with typed transitions from
$\Sigma$, an outcome class $\theta \in \Theta$, and a cycle depth $D(\sigma)$, such that every
walk step is grounded in a canonical event of Equation~\eqref{eq:kinds} and every terminal
class is decided by an explicit rule over terminal events, tool-error counts, or (for
evaluation dialects) environment-verified outcomes.
\end{definition}

\subsection{Outcome classes and cycle depth}
The walk ends in exactly one terminal class.
Contained ($\theta_{\mathrm{cont}}$) requires an explicit close with no unresolved tool
errors; handoff ($\theta_{\mathrm{hand}}$) is forced by handoff-typed tools or events;
abandoned ($\theta_{\mathrm{aban}}$) is the absence of a terminal event at stream end;
compliance-escape ($\theta_{\mathrm{esc}}$) is the audit overlay inherited from the companion
grammar \citep{garcia2026grammar}, computed as a post-hoc check rather than a walk state.
Cycle depth $D(\sigma)$ counts regressions to settled states --- the drift statistic that
distinguishes a healthy clarify loop from an endless one.

\subsection{The emission-model assumption, stated honestly}
The guarantees of the companion grammar attach to its \emph{stopping} layer under calibrated
emissions \citep{garcia2026grammar}.
$\tau$ supplies deterministic emissions (pattern tables), so this paper's coverage results
inherit the walk \emph{semantics} only --- state identity, transition typing, outcome
decidability --- and none of the prediction guarantees.
This is the honest dependency direction: the coverage metric is defined on walks, its
correctness claim is invariant to how walks are produced, and a calibrated state estimator
can replace \texttt{infer\_states()} without changing any downstream definition.
What the baseline does add is reproducibility: the same export always yields the same walk,
which Section~\ref{sec:evaluation} exploits to make sensitivity sweeps exact.
